\section{Frontera de investigación: de agentes web a agentes ejecutores multimodales generales}

\subsection{Abstracciones de acción de nivel más alto}
El ponente propone una dirección que merece una inversión a largo plazo: diseñar para los agentes abstracciones semánticas de nivel más alto, de modo que el modelo no tenga que buscar desde acciones atómicas de click/type en cada tarea. Esto puede reducir de manera considerable la profundidad de búsqueda y la complejidad de ingeniería.

\begin{importantbox}{Beneficios potenciales de las abstracciones de nivel alto}
\begin{itemize}
    \item Reducen el ruido de las acciones de bajo nivel;
    \item Mejoran la generalización entre sitios;
    \item Crean una interfaz unificada para la colaboración entre múltiples agentes.
\end{itemize}
\end{importantbox}

\subsection{Búsqueda en paralelo y ejecución en múltiples instancias}
El ponente también mencionó la importancia de la paralelización: la misma tarea puede explorarse en múltiples instancias siguiendo rutas distintas, para después devolver la trayectoria óptima al flujo principal. Esto coincide con la tendencia de crecimiento del ``test-time compute'' en la era de los modelos grandes.

\begin{knowledgebox}{Puntos clave del diseño de la ejecución en paralelo}
\begin{enumerate}
    \item Las subtareas pueden reproducirse de forma independiente;
    \item Se comparten señales de evaluación intermedias;
    \item Se unifica la estrategia de resolución de conflictos (por ejemplo, tomar la ruta de mayor valor).
\end{enumerate}
\end{knowledgebox}

\subsection{El siguiente paso de los benchmarks de evaluación}
La próxima generación de benchmarks necesita cubrir de manera más realista: sitios web dinámicos, páginas personalizadas, sesiones de larga duración, interacción entre dispositivos y colaboración con múltiples roles. De lo contrario, el sistema mejorará en benchmarks estáticos pero fallará en el entorno de producción real.

\begin{warningbox}{La brecha estructural entre los benchmarks y el despliegue real}
Los productos web cambian de manera continua, los elementos se cargan de forma dinámica y las restricciones de permisos y sesión son complejas; cualquier benchmark fijo puede quedar obsoleto con rapidez. El sistema de evaluación debe admitir actualización continua y seguimiento de versiones.
\end{warningbox}

\subsection{Resumen de la sección}
El punto de avance a largo plazo de los agentes multimodales está en la abstracción de acciones, la búsqueda en paralelo y el sistema de evaluación dinámica. La competencia futura ocurrirá cada vez más en el nivel de la ingeniería de sistemas, y no solo en el nivel de la escala de parámetros.
